\section{十一固定点と五対の最大monodromy}
\label{sec:master-monodromy-v25}

前節の標準形から、指定反対合を持つ枠付き関係族上で $Z_2$ の点monodromyを決定する。
粗モジュライ上の普遍被覆の存在は仮定しない。

\subsection{重みだけを動かす十一点被覆}

単純正規交叉をなす滑らかな二次曲線 $Q=V(q)$ と五本の直線
$H_i=V(\ell_i)$ を固定する。$U=\mathbb P^2\setminus(Q+\sum H_i)$ と置き、
重みを射影座標 $[w_1:\cdots:w_5]\in\mathbb P^4$ として
\[
 \theta_i=d\log\ell_i-\frac12d\log q,\qquad
 \omega_w=\sum_iw_i\theta_i
\]
と書く。$\sum_iw_i\ne0$ の開部分では、その和を一に正規化して前節の表示に戻る。

\begin{lemma}[臨界点対応の既約性と次数]
\label{lem:master-transitive-v25}
この固定した配置について、一般の重みの臨界点は相異なる十一点である。
それらの幾何学的monodromyは推移的である。
\end{lemma}
\begin{proof}
最初の三本の直線の方程式を射影座標に取れる。
$U$ 上では
\[
 \theta_1-\theta_3=d\log(\ell_1/\ell_3),\qquad
 \theta_2-\theta_3=d\log(\ell_2/\ell_3)
\]
が余接空間の基底を与える。従って評価写像
$\mathcal O_U^5\to\Omega_U^1$ は全射で、その核は階数三の束である。
臨界点対応
\[
 \mathcal C=\{(x,[w])\in U\times\mathbb P^4:\omega_w(x)=0\}
\]
はこの核の射影束なので滑らかかつ既約である。

$D_{\rm pol}=Q+\sum_iH_i$ とする。$\omega_w$ は
$E=\Omega_{\mathbb P^2}^1(\log D_{\rm pol})$ の切断である。
各直線の留数は $w_i$、二次曲線の留数は $-\frac12\sum_iw_i$ なので、
これらが全て非零なら零点は境界と交わらない。
単純正規交叉の対数束のChern類から
\[
 c\bigl(T_{\mathbb P^2}(-\log D_{\rm pol})\bigr)
 =\frac{1+3H+3H^2}{(1+2H)(1+H)^5}
 =1-4H+11H^2.
\]
従って $c_2(E)=11H^2$ である。これが非零なので切断の零点は空ではなく、
境界の留数が非零な重みの開部分へ $\mathcal C$ は全射である。
特に射影は支配的である。
評価写像の全射性と標数零の一般滑らかさにより、
一般の切断の零点は有限被約となり、その長さは十一である。
境界を避ける開部分で零スキームは射影平面の閉部分族なので固有であり、
さらに重みの開部分を縮めれば有限エタール被覆になる。
$\mathcal C$ の既約性からこの被覆は連結であり、monodromyは推移的である。
\end{proof}

この次数は $\chi(U)=11$ と一致する。
一般のmaster functionの臨界点計数は既知の理論であり
\cite[Theorem~1]{HuhML}、ここでは後の群論に必要な既約性と次数を直接確認した。

\begin{lemma}[単独の二重点を持つ有理数例]
\label{lem:master-single-fold-v25}
アフィン座標 $x=X/Z,y=Y/Z$ で $q=1+x^2+y^2$ とする。
\[
 \ell_i=1+a_ix+b_iy,\qquad
 ((a_i,b_i))_{i=1}^5=((3,-1),(2,-3),(-1,-3),(-3,1),(3,2))
\]
および
\[
 (w_i)=\left(-\frac{53}{16},\frac{17}{6},-\frac{25}{48},
                    \frac38,\frac{13}{8}\right)
\]
を取る。この配置は単純正規交叉であり、$\sum_iw_i=1$ である。
$\omega_w$ の零スキームは原点で長さ二、他の九点で被約であり、
全ての点は $U$ に含まれる。
重み方向 $(1,-1,0,0,0)$ はこの二重点を単純なfoldとして開く。
\end{lemma}
\begin{proof}
原点で $\omega_w=0$ であり、局所的な対数関数のHessianと、
その核方向の三階微分は
\[
 \operatorname{Hess}_0=
 \begin{pmatrix}0&0\\0&-203/8\end{pmatrix},
 \qquad \partial_x^3\!\left(\sum_iw_i\log\ell_i-\tfrac12\log q\right)_0=-65.
\]
重み方向 $(1,-1,0,0,0)$ における $dx$ 成分の定数項の微分は
$a_1-a_2=1$ である。
従って $y$ を陰関数定理で消すと、残る方程式は
非零の二次項を持ち、この重み方向でその定数項が横断的に変化する。
これは複素解析的な単純foldである。

他に重複点がないことと、無限遠を含む全零点の確認は、
次の小さい有理数証明書で行う。
$P=\prod_i\ell_i$ と置き、
\[
 f_x=q\sum_iw_i a_i P/\ell_i-xP,\qquad
 f_y=q\sum_iw_i b_i P/\ell_i-yP
\]
を取る。$\mathbb Q[t,y,x]$ 上で
$(f_x,f_y,tqP-1)$ の辞書式Gr\"obner基底を計算すると、
$t$ を含まない部分は
\[
 y-r(x),\qquad x^2h_9(x)
\]
という形である。ここで $\deg r=10$、$r(0)=r'(0)=0$、
$h_9$ は九次のモニック多項式で
\[
 h_9(0)=-\frac{1459976}{1929737}\ne0.
\]
その係数は $101$ で整であり、剰余は
\[
 \bar h_9=x^9+34x^8-27x^7-46x^6+13x^5-42x^4
                    +19x^3-28x^2-40x-30.
\]
$\operatorname{disc}(\bar h_9)=18\ne0$ なので $h_9$ は分離的である。
飽和により $qP$ はこの商で単元であり、アフィン零スキームは
原点の長さ二と九つの相異なる点に分かれる。

無限遠は同次化した一形式の三分子で検査する。
$X=1,Z=0$ で残る二分子の最大公約数から境界の因子を除くと定数となり、
残る点 $[0:1:0]$ でも二分子は同時に零でない。
境界上では全ての留数が非零である。
従ってこれが射影平面上の全零スキームである。
検証器 \nolinkurl{master_monodromy_v25.py} と付属のJSONは、
有理数上の全Gr\"obner基底、無限遠の最大公約数、
配置の単純正規交叉性、上記の微分を再計算する。
\end{proof}

\begin{proposition}[固定十一点の完全対称群]
\label{prop:master-S11-v25}
補題\ref{lem:master-single-fold-v25}の二次曲線と五直線を固定し、
重みだけを変化させた十一点被覆の幾何学的monodromyは $\mathfrak S_{11}$ である。
このmonodromyは、二次曲線上の十個の非固定点を一つずつ固定する。
\end{proposition}
\begin{proof}
補題\ref{lem:master-transitive-v25}から群は十一点上推移的である。
次数十一は素数なので、この作用は原始的である。
補題\ref{lem:master-single-fold-v25}の単独foldの周囲の小さなループは、
原点から分裂する二点を交換し、他の九点を固定する。
従って群は一つの互換を含む。

原始的置換群が互換を含めば全対称群であることも直接示せる。
その互換の全共役を辺とするグラフを十一個の頂点上に作る。
連結成分は群によって置換されるblockをなし、辺があるので原始性からグラフは連結である。
連結グラフの辺の互換は全対称群を生成する。

このループでは $q,\ell_1,\ldots,\ell_5$ は固定している。
命題\ref{prop:master-critical-fixed-v25}より十個の非固定点は
$Q\cap\bigcup_iH_i$ なので、それぞれ固定される。
\end{proof}

\subsection{五対の独立な交換と全群}

\begin{theorem}[$Z_2$ の最大monodromy]
\label{thm:master-full-monodromy-v25}
$Z_2$ の被約二十一点開部分における、指定反対合と可換な点monodromyは
\[
 \boxed{\;
 G_{Z_2}=\mathfrak S_{11}\times
       \bigl((\mathbb Z/2\mathbb Z)^5\rtimes\mathfrak S_5\bigr).
 \;}
\]
第一因子は十一固定点に作用し、第二因子は二次曲線上の五対に作用する。
\end{theorem}
\begin{proof}
前節の構成と逆構成は同じ成分 $Z_2$ の非空開部分で成立する。
補題\ref{lem:master-single-fold-v25}の有理数例でも、
自己反対incidence、$\operatorname{rank}F=39$、三次関係写像の階数三十六、
復元写像の階数四十八を検証している。
従って、その重み変形はこの成分の実際の点被覆へ写る。
命題\ref{prop:master-S11-v25}により
$\mathfrak S_{11}\times1$ は $G_{Z_2}$ に含まれる。

次に $Q$ を固定して五直線を動かす。
一本の直線を双対二次曲線の滑らかな点の周囲で一周させると、
その直線と $Q$ の二交点が交換される。
他の四直線を固定してこの操作を独立に行えるので、
五対への作用は $(\mathbb Z/2\mathbb Z)^5$ を含む。
五つの $(H_i,w_i)$ を同時に置換するbraidにより、
対の置換群 $\mathfrak S_5$ も得る。

枠付きの族上で最後のループを閉じる点を明記する。
ラベルの置換は負固有空間の行基底の置換と、関係空間 $K$ の基底変更で実現される。
これらは指定反対合を保つ連結な一般線形群に属する。
その群内の道によって終点を始点へ戻せば、点配置を一つずつ固定する区間を付け加えて
枠付きの族でもループが閉じる。
必要な階数条件を満たさない閉部分は除外できる。
一般の複素代数多様体で非空Zariski開部分へ制限しても、
既存の有限エタール被覆のmonodromy像は小さくならない。
従って第二因子への射影は全wreath productに全射である。

もとの上界は
$\mathfrak S_{11}\times((\mathbb Z/2\mathbb Z)^5\rtimes\mathfrak S_5)$ である。
その部分群が第一因子全体を含み、第二因子へ全射なら、
第一成分を消すことにより各 $(1,g)$ も含む。
よって直積全体と一致する。
\end{proof}

この証明は有理数上の特殊化Galois群から複素monodromyを推定するものではない。
連結な臨界点対応、複素解析的な単独fold、直線と二次曲線の交点被覆を使った
幾何学的な証明である。$Z_1,Z_3$ のmonodromyについて等号を主張するものでもない。
